
     \documentclass[fleqn]{article}      
      \usepackage{amsmath}
      \usepackage{fontspec}
      \usepackage{kotex} % 한국어 지원  

      \begin{document}      
        

쉬움: \\\\  
1. 함수는 두 집합 사이의 관계로, 각 입력값에 대해 고유한 출력값이 대응되는 규칙이다. \\\\  
   해설: 함수는 일반적으로 $f : A \to B$ 형태로 표현되며, $A$는 정의역, $B$는 공역이다. \\\\  
2. 함수의 그래프는 각 입력에 대해 하나의 출력만 가지므로, 원형 그래프는 함수가 아니다. \\\\  
   해설: 모든 입력에 대해 고유한 출력이 있어야 함수의 조건을 충족한다. \\\\  
3. 가능한 조합의 수는 $3^3 = 27$이다. \\\\  
   해설: 각 입력값에 대해 3개의 선택이 가능하므로, 총 경우의 수는 $3^3$이다. \\\\  
.......... \\\\  

중간: \\\\  
1. 기울기는 2이다. \\\\  
   해설: 함수의 기울기는 $2$로, 이는 $x$가 1 증가할 때 $y$가 2 증가함을 의미한다. \\\\  
2. $f(1) = 1^2 - 4 = -3$이다. \\\\  
   해설: 주어진 함수에 $x = 1$을 대입하여 계산하였다. \\\\  
3. 정의역은 $x \neq 0$이다. \\\\  
   해설: $g(x) = \frac{1}{x}$에서 분모가 0이 되지 않는 실수의 집합이다. \\\\  
.......... \\\\  

어려움: \\\\  
1. 극대값은 $f(1) = 2$, 극소값은 $f(-1) = -4$이다. \\\\  
   해설: 미분을 통해 극값을 찾고, 대입하여 극대와 극소를 확인하였다. \\\\  
2. $h(0) = 2^0 = 1$이다. \\\\  
   해설: 지수함수의 기본 성질에 의해 $2^0 = 1$이다. \\\\  
3. 교점은 두 함수가 같아야 하므로, $3x - 4 = 2x + 2$를 풀면 $x = 6$이다. \\\\  
   해설: 두 함수를 같게 놓고 $x$의 값을 구하여 교점을 찾았다. \\\\  
.......... \\\\        
      \end{document} 
  